\section{Implementation}
\label{sec:method}

All three repairs live in \code{ideas.py} and are built from the documented update pipeline
(\code{\_encode $\to$ \_evaluate $\to$ \_votes $\to$ \_select\_feedback $\to$
\_feedback\_counts $\to$ \_commit}) and \code{functional.apply\_feedback}. No library file is
modified. The architecture is exactly the one evaluated in the first report --- layer~1
$\pOne\times\pOne$ patches with \cOne\ shared clauses, $\poolK\times\poolK$ OR pool, layer~2
$\pTwo\times\pTwo$ patches with \cTwo\ clauses --- so every number here is directly comparable
to its E1 counterpart.

\subsection{\ideaA\ --- a patch autoencoder for layer 1}
\label{sec:method-auto}

Each patch is split into a \emph{target} --- its centre pixels, across all bit planes --- and
a \emph{context}, the surrounding ring. A flat coalesced Tsetlin machine with one output per
target bit is trained multi-label to predict the target from the context, and its clauses are
then written into the convolutional layer. No label is used anywhere.

Two choices are worth stating. The split is \emph{spatial} rather than a random subset of
bits, because thermometer planes of the same pixel are nested ($v>t_1$ implies $v>t_0$):
holding out one plane of a pixel while showing another leaks the answer, and the clause learns
the encoder instead of the image. And the trained clauses are embedded into a layer whose
literal space is the \emph{full} patch plus position bits, identical to a greedily trained
layer~1; the held-out pixels and the position bits simply stay excluded. The two layer-1
variants are therefore the same model with the same capacity, differing only in objective.

\subsection{\ideaB\ --- a two-sided density controller}
\label{sec:method-calib}

A clause fires on a patch iff none of its included literals is false there, so its firing rate
is monotone in its literal set: removing a literal can only raise it, adding one can only lower
it. The target rate is therefore reachable by a greedy walk in \emph{both} directions, which is
what ``adjust each clause towards it'' asks for. Per round, for the clauses below the band,
\[
  \mathrm{gain}[j,k] \;=\; \#\{\text{patches where clause $j$ has exactly one violated literal,
  and it is $k$}\}
\]
names the literal whose removal buys the most; it is the matmul $(\text{violations}=1)^{\!\top}
(1-\text{literals})$ masked to included literals, the same shape of computation as the Type~Ia
count in \code{functional}. Where no single removal helps, the most-often-violated literal goes
instead, so a round always makes progress. For clauses above the band, including literal $k$
drops the rate by the fraction of currently-matching patches in which $k$ is false, so the
literal landing closest to the target without undershooting is added.

The same loop takes a second target. A clause's firing rate and its \emph{size} are different
properties --- a long conjunction that is rare-but-common and a short general one can sit at
the same rate --- so the controller accepts a size cap alongside the rate, and
\S\ref{sec:sizetarget} uses it to test which of the two the second layer actually needs.
Under a cap the removal rule changes: dropping the most load-bearing literal leaves the most
permissive ones and sends the rate towards~1, dropping the least leaves the most selective and
sends it towards~0, so the literal that leaves the rate closest to target goes instead. That
choice is exact rather than heuristic, because removing literal $k$ lets through precisely the
patches it alone was blocking.

Two details decide whether this is a controller or a gimmick. \lead{Added literals are
restricted to patch pixels.} The cheapest way to hit any firing rate is a position literal:
\code{y > 27} fires on exactly $1/29$ of patch rows whatever the image contains. Left free, the
controller takes it every time and produces channels that carry no information at all.
\lead{A clause that has once been inside the band is frozen}, so the two rules cannot chase
each other. Adjustments are written as ``included'' / ``excluded, one step below the
boundary'' --- exactly where a fresh automaton sits --- so calibration stays inside the
automata state and survives continued training, which is what lets it be interleaved with
\ideaC.

\subsection{\ideaC\ --- credit propagation}
\label{sec:method-credit}

If layer-2 clause $j$ matched patch $(p_y,p_x)$ and includes the positive literal ``channel $k$
fires at offset $(dy,dx)$'', the reason it matched is that layer-1 clause $k$ fired somewhere
inside the pooled cell $(p_y{+}dy,\,p_x{+}dx)$. Whatever feedback layer~2 decided clause $j$
deserves is therefore also a verdict on layer-1 clause $k$ at that position. The rule maps each
of the three Tsetlin feedback types down through the same inclusion link:

\begin{table}[h]
\centering\small
\caption{The credit rule. Each layer-2 event resolves to a single layer-1 position by three
draws: a matching layer-2 patch, one included positive channel literal of that clause, and one
layer-1 position inside the pooled cell that literal names at which the channel actually fired.}
\label{tab:creditrule}
\begin{tabular}{@{}lll@{}}
\ttoprule
\thead{layer-2 event} & \thead{layer-1 clause $k$ receives} & \thead{why} \\
\ttmidrule
Type Ia (matched)             & Type Ia at that position & reinforce the pattern that fired \\
Type Ib (should have matched) & Type Ib                  & a channel it needs did not fire \\
Type II (matched, wrong class)& Type II at that position & make it stop firing there \\
\ttbotrule
\end{tabular}
\end{table}

Four choices, each of which the experiments turned out to depend on.

\lead{One credited literal per event.} A layer-2 clause can include hundreds of literals;
crediting all of them would make a single event a bulk update of layer~1. One included
positive literal is drawn uniformly, exactly as the convolutional machine draws one matching
patch per event.

\lead{Type Ib is blamed, not sampled.} A Type Ib event says the layer-2 clause \emph{should}
have matched and did not, so the credit question is which literal blocked it. The clause is
taken at the patch it came closest to matching (fewest violated literals) and the blame falls
on an included positive channel literal that was \emph{false} there. Crediting a random
included literal instead --- including ones that were satisfied --- blames layer-1 clauses
that did their job, and erases layer~1 within an epoch (\S\ref{sec:credit}).

\lead{One event per (example, layer-1 clause).} In the classical algorithm an example gives a
clause at most one feedback event. Without this, \cTwo\ layer-2 clauses credit \cOne\ layer-1
clauses on every image and layer~1 sees orders of magnitude more events than it was designed
for.

\lead{A step size.} Layer-2 clauses are conjunctions over \emph{named} channels. When layer-1
clause $k$ changes, every layer-2 clause that includes channel $k$ silently becomes a different
rule. The fraction of credit events actually applied is the only knob that trades layer-1
progress against that drift, so it is swept rather than assumed.
